\documentclass[11pt,a4paper]{article}
\usepackage{xeCJK}
\usepackage{fontspec}
\usepackage{xcolor}
\usepackage{geometry}
\usepackage[protrusion=true]{microtype}
\geometry{margin=1.8cm,top=1.9cm}
\linespread{1.25}

\setmainfont{Baskerville}
\setCJKmainfont{Songti SC}

\definecolor{tipblue}{RGB}{0,95,160}
\definecolor{rulegray}{RGB}{190,190,190}

\setlength{\parindent}{0pt}
\setlength{\parskip}{0.35em}
\newcommand{\num}[1]{\makebox[1.9em][r]{#1.}\hspace{0.3em}}
\newcommand{\wline}{\noindent\rule{\linewidth}{0.4pt}}
% 中文题干：宋体粗体；提示词：蓝色小字；英文句子：Baskerville 正体
\newcommand{\prompt}[1]{\emph{\textcolor{tipblue}{\small（#1）}}}

\begin{document}
\pagestyle{plain}
\begin{center}
{\LARGE\bfseries 上海高考高分翻译背诵 \quad 57}\\[0.25em]
{\large 每日背诵-57　2026-08-17 \quad\ 默写（中 $\to$ 英）}\\[0.35em]
{\footnotesize 来源：微信公众号「发现之旅 速来记单词」· 2026-08-17 · 赵文瀚}
\end{center}

\vspace{0.4em}
\noindent\begin{tabular}{@{}p{0.33\textwidth}@{}p{0.34\textwidth}@{}p{0.33\textwidth}@{}}
姓名：\underline{\hspace{2.0cm}} & 日期：\underline{\hspace{2.0cm}} & 得分：\underline{\hspace{1.4cm}}\,/\enspace 5 \\
\end{tabular}

\vspace{0.5em}
\noindent{\small 根据中文句子与提示词写出英文翻译，题号沿用原文。}

\vspace{0.8em}
\num{1}\textbf{不过矛盾还是明显存在。}

\vspace{0.3em}

\wline\vspace{0.3em}
\wline

\num{2}\textbf{他并不隶属于任何大学。}

\vspace{0.3em}

\wline\vspace{0.3em}
\wline

\num{3}\textbf{当然，对往事的怀念也是原因之一。}

\vspace{0.3em}

\wline\vspace{0.3em}
\wline

\num{4}\textbf{这座跨海大桥在落日的映衬下多么壮丽啊！}\prompt{against}

\vspace{0.3em}

\wline\vspace{0.3em}
\wline

\num{5}\textbf{在失重的环境中，吃饭喝水皆非易事，但中国人对吃的挚爱让我们的宇航员在空间站里成功地吃上了烧烤盛宴。}\prompt{so…that…}

\vspace{0.3em}

\wline\vspace{0.3em}
\wline

\end{document}
